\addcontentsline{toc}{subsubsection}{Solution 31. The Iterated Parallelogram--Quadratic Squeeze}

\begin{dlsummary}[title={\bfseries Solution 31. The Iterated Parallelogram--Quadratic Squeeze},breakable]

\begin{proof}
\textbf{Step 1. Derive a quadratic lower bound.}

Comparing the outer expressions in the given chain gives
\[
f(x)-x^2+f(f(x))-f(y)+2y^2-2yf(x)\geq0.
\]
Setting $y=f(x)$ cancels the terms involving $f(f(x))$ and the quadratic
terms in $y$, leaving
\[
f(x)\geq x^2\qquad(x\in\mathbb R). \tag{1}
\]
Define $g(x)=f(x)-x^2$. Then $g(x)\geq0$ for every real $x$.

\medskip
\noindent
\textbf{Step 2. Obtain midpoint concavity of the residual.}

The right-hand inequality and the identity
\[
\frac{(x+y)^2+(x-y)^2}{2}=x^2+y^2
\]
show that
\[
g(x+y)+g(x-y)\leq2g(x)\qquad(x,y\in\mathbb R). \tag{2}
\]

\medskip
\noindent
\textbf{Step 3. Prove that $g$ is constant.}

Fix $a,b\in\mathbb R$, and write
\[
x_k=a+k(b-a)\quad(k\geq0),
\qquad
\Delta_k=g(x_{k+1})-g(x_k).
\]
For each $k\geq1$, apply (2) with $x=x_k$ and $y=b-a$. This gives
\[
g(x_{k+1})+g(x_{k-1})\leq2g(x_k),
\]
so $\Delta_k\leq\Delta_{k-1}$. Consequently,
$\Delta_k\leq\Delta_0=g(b)-g(a)$ for all $k\geq0$. For every positive
integer $n$, nonnegativity of $g$ now yields
\[
-g(a)\leq g(x_n)-g(a)
=\sum_{k=0}^{n-1}\Delta_k
\leq n\bigl(g(b)-g(a)\bigr).
\]
Thus $g(b)-g(a)\geq-g(a)/n$ for every $n$, and letting $n\to\infty$
shows $g(b)\geq g(a)$. Since $a,b$ were arbitrary, interchanging them
proves $g(a)=g(b)$. Therefore $g(x)=d$ for a constant $d\geq0$, and
\[
f(x)=x^2+d. \tag{3}
\]

\medskip
\noindent
\textbf{Step 4. Use the composition condition and verify.}

By (3),
\[
f(f(2))=(4+d)^2+d=16+9d+d^2\leq16.
\]
Since $d\geq0$, this forces $d=0$. Hence $f(x)=x^2$ for every real $x$.

For this function, $f(f(2))=16$, the middle and right expressions in
the given chain are both $x^2+y^2$, and the left expression equals
\[
x^2+y^2-(y-x^2)^2\leq x^2+y^2.
\]
Thus $f(x)=x^2$ satisfies all the conditions, completing the proof.
\end{proof}

\end{dlsummary}